\documentclass[14pt,a4paper]{extarticle}
\usepackage[T2A]{fontenc}
\usepackage[utf8]{inputenc}
\usepackage[english,russian]{babel}
\usepackage{lmodern}
\usepackage{geometry}
\usepackage{array}
\usepackage{tabularx}
\usepackage{longtable}
\usepackage{ltablex}
\usepackage{booktabs}
\usepackage{hyperref}
\usepackage{indentfirst}
\usepackage{microtype}
\usepackage{graphicx}
\usepackage{setspace}
\usepackage{listings}
\usepackage{graphicx}
\usepackage{caption}
\usepackage{float}
\usepackage{tocloft}
\usepackage{titlesec}
\usepackage{newtxtext,newtxmath}

\titleformat{\section}        % форматируем \section
  {\normalfont\bfseries}      % стиль: обычный + жирный
  {\thesection.}              % номер секции с точкой
  {2mm}                       % расстояние между номером и текстом
  {}                           % перед заголовком можно добавить что-то
\titlespacing*{\section}{1.5cm}{*3}{*1}  % {отступ слева}{отступ сверху}{отступ снизу}

\captionsetup{justification=centering}
\renewcommand{\ttdefault}{cmtt}
\keepXColumns
\setlength{\parindent}{3em}
\newgeometry{left=2.5cm,right=2.5cm,top=3cm,bottom=2.5cm}

\renewcommand{\contentsname}{\hspace*{4em}{\smallСодержание}}
\setlength{\cftbeforesecskip}{0.5em}
\renewcommand{\cftsecleader}{\cftdotfill{\cftdotsep}}
\renewcommand{\cftsecaftersnum}{.\hspace{2mm}}

\title{Отчёт по дисциплине "Основы предпринимательства"}
\author{Выполнил студент группы М-РИТ-25-1, Станиславчук Сергей}

\begin{document}

\setstretch{1.5}
\noindent
\begin{minipage}[t]{0.01\textwidth}
    \includegraphics{~/programmer/++Programmer/5_1/LSTU.jpg}
\end{minipage}%
\begin{minipage}[c]{1.1\linewidth}
    \centering
    {\small
    МИНИСТЕРСТВО НАУКИ И ВЫСШЕГО ОБРАЗОВАНИЯ РФ \\[1em]
    ФЕДЕРАЛЬНОЕ ГОСУДАРСТВЕННОЕ БЮДЖЕТНОЕ \\
    ОБРАЗОВАТЕЛЬНОЕ УЧРЕЖДЕНИЕ ВЫСШЕГО ОБРАЗОВАНИЯ \\
    «ЛИПЕЦКИЙ ГОСУДАРСТВЕННЫЙ ТЕХНИЧЕСКИЙ УНИВЕРСИТЕТ»
    }
\end{minipage}%

\vspace{3em}

\noindent
\hspace{1em}%
\begin{minipage}[t]{0.8\textwidth}
\textbf{Институт} \hspace{2em} компьютерных наук \\
\textbf{Кафедра} \hspace{2em} автоматизированных систем управления
\end{minipage}


\vspace{6em}

\begin{center}
\textbf{ПРАКТИЧЕСКАЯ РАБОТА №7} \\
по дисциплине \\ <<Основы предпринимательства>>
\end{center}

\vspace{6em}

\noindent
\begin{tabular*}{\textwidth}{@{\extracolsep{\fill}} l c r }
    Студент М-РИТ-25-1 & \makebox[4cm]{\hrulefill} & Станиславчук С.М. \\[2em]
Руководитель & \\
Канд. юр. наук, доцент & \makebox[4cm]{\hrulefill} & Моисеева И.И. \\
\end{tabular*}

\vfill
\begin{center}
    Липецк, 2025
\end{center}

\thispagestyle{empty}

\newpage
\tableofcontents
\newpage

\section{Задание 1}
\begin{sloppypar}
    При вложении денежных средств в мероприятие А из 150 случаев прибыль в
    сумме 20,0 тыс. руб. была получена в 75 случаях, прибыль 25,0 тыс. руб. в 60 случаях и
    прибыль 30,0 тыс. руб. в 15 случаях. Найдите среднеожидаемое значение прибыли.
\end{sloppypar}

Решение.

Среднеожидаемое значение прибыли = 20.0 * 0.5 + 25.0 * 0.4 + 30.0 * 0.1 = 23

    Ответ: 23 тыс. руб

\section{Задание 2}
\begin{sloppypar}
    Осуществление мероприятия Б из 150 случаев давало прибыль 19,0 тыс, руб. в 60
    случаях, прибыль 24,0 тыс, руб. — в 45 случаях, 31,0 тыс, руб. в 45 случаях. Какова
    ожидаемая прибыль? Сравните величины ожидаемой прибыли при вложении средств в
мероприятие А и Б.
\end{sloppypar}

Решение.

Математическое ожидаение для мероприятия Б = (19 * 60 + 24 * 45 + 31 * 45) / 150 = 24.1

Вывод: вепичина получаемой прибыли при мероприятии А колеблется от 20,0 до 30,0 тыс. руб., средняя величина составляет 23 тыс. руб.; в мероприятии Б величина получаемой прибыли колеблется от 19,0 до 31,0 тыс. руб. и средняя величина равна 24,1 тыс. руб.


\section{Задание 3}
\begin{sloppypar}
    Организация намерена инвестировать часть капитала в ценные бумаги А и В, при
    этом ожидаемая доходность ценной бумаги А составляет 7\%, а среднеквадратическое
    отклонение доходности составляет 5\%; ожидаемая доходность ценной бумаги В — 12
    \%, а среднеквадратическое отклонение доходности — 8\%. Какова ожидаемая прибыль?
\end{sloppypar}

Решение.

Ожидаемая прибыль = 0.5 * 0.07 + 0.5 * 0.12 = 9.5\%

Ответ: 9.5\%


\section{Задание 4}
\begin{sloppypar}
    При вложении денежных средств в мероприятие А из 200 случаев прибыль 25
    тыс. руб. была получена в 20 случаях, прибыль 30 тыс. руб. была получена в 80 случаях,
    прибыль 40 тыс. руб. была получена в 100 случаях. При вложении капитала в
    мероприятия Б из 240 случаев прибыль 30 тыс. руб. была получена в 144 случаях,
    прибыль 35 тыс. руб. была получено в 72 случаях, прибыль 45 тыс. руб. была получена в
    24 случаях. Выберите вариант вложения капитала в мероприятие. Критерием выбора
    является наибольшая сумма средней прибыли.
\end{sloppypar}

Решение.

Мероприятие A = (25 * 20 + 30 * 80 + 40 * 100) / 200 = 34.5
Мероприятие Б = (30 * 144 + 35 * 72 + 45 * 24) / 240 = 33.0

Выбор однозначен — мероприятие А, так как средняя прибыль выше.

Ответ: мероприятие A.


\section{Задание 5}
\begin{sloppypar}
    Рассчитайте коэффициент риска и выбрать наименее рисковый вариант вложения
    капитала, если известно, что по варианту А собственные средства инвестора составляют
    5 000 тыс, руб., возможная сумма убытка составляет 3 500 тыс, руб.; по варианту Б —
    собственные средства 7 500 тыс, руб., возможная сумма убытка 4 000 тыс, руб.
\end{sloppypar}

Решение.

Коэффициент риска считаем как отношение возможного убытка к собственным средствам:

\[
K = L / S
\]

Вариант A = (3500) / 5000 = 0.70
Вариант Б = (4000) / 7500 = 0.53

Наименее рисковый вариант — Б, так как коэффициент риска меньше.

Ответ: Вариант Б.

\section{Задание 6}
\begin{sloppypar}
    Предприятие выпускает технически сложную продукцию, производство которой
    требует значительного количества комплектующих деталей. Предприятие может
    покупать эти детали по цене 100 тыс. руб./шт. или производить их у себя.
    Определите, что более выгодно для предприятия (порог рентабельности
    пройден). Себестоимость производства деталей на самом предприятии: прямые
    переменные издержки 85 тыс. руб./шт.; постоянные накладные расходы 20 тыс.
    руб./шт.; итого полная себестоимость 105 тыс. руб./шт.
\end{sloppypar}

Решение.

Цена покупки = 100 т.руб
Производство = 85 т.руб
Разница = 100 - 85 = 15 т.руб

Производить детали на собственном предприятии выгоднее — экономия 15 тыс. руб. на каждую деталь, несмотря на то, что полная себестоимость (105 тыс. руб.) выше цены покупки.


\section{Задание 7}
\begin{sloppypar}
Определите тип финансовой устойчивости предприятия, рассчитан
коэффициенты: автономии, маневренности собственных оборотных средств,
соотношения заемных и собственных средств, обеспеченности собственными
оборотными средствами, на основании исходных данных, сделать выводы (табл. 14)
\end{sloppypar}

Актив / тыс. руб. / Пассив / Тыс. руб
1. Основные редства и нематериальные активы / 8 000 / 1. Собственный капитал / 10 000
2. Запасы и затраты / 5 000 / 2. Долгосрочные кредиты / 1 500
3. Дебиторская задолженность / 3 000 / 3. краткосрочные кредиты / 1000
4. Краткосрочные финансовые вложения / 1 500 / 4.5. Кредиторская задолженность / 6000
5. Денежные средства / 1000 /
Баланс / 18 500 / Баланс 18 500

Решение.

Собственные оборотные средства COC
COC = CK - BA = 10000 - 8000 = 2000

Коэффициент автономии K = CK / Баланс = 10000 / 18 500 = 0.54

Коэффициент маневренности M = COC / CK = 2000 / 10000 = 0.20

Коэффициент соотношения заемных и собсвтенных средств S = ЗК / CK = 8 500 / 10 000 = 0.85

Коэффициент обеспеченности собственными оборотными средствами O = COC / OA = 2000 / 10 500 = 0.19

Сравнение :  COC = 2000, запасы = 5000. COC < 3

Вывод:

Собственных оборотных средств недостаточно, но предприятие покрывает запасы за счёт заемных источников.

Тип финансовой устойчивости:

Нормальная финансовая устойчивость


\section{Задание 8}
\begin{sloppypar}
Оцените финансовый риск предприятия с помощью коэффициентов абсолютной и
текущей ликвидности, если известно, что денежные средства предприятия составляют 3
тыс, руб., стоимость ценных бумаг —6 тыс, руб., дебиторская задолженность — 3 500
руб., запасы на предприятии составляют 1 800 руб., кредиторская задолженность — 3
600 руб., краткосрочные займы предприятия — 6 тыс, руб. Сделать выводы.
\end{sloppypar}

Решение.

KO = 9.6

OA = 14.3

Коэффициент абсолютной ликвидности K = (ДС + ЦБ) / KO = (3 + 6) / 9.6 = 0.94

Значение резко выше нормы — предприятие легко покрывает краткосрочные обязательства мгновенно.

Коэффцицент текущей ликвидности T = OA / KO = 14.3 / 9.6 = 1.49

Значение ниже нормы.

Вывод: Предприятие финансово устойчиво, угрозы неплатёжеспособности нет.



\section{Задание 9}
\begin{sloppypar}
Рассчитать риск банкротства предприятия при условии, что собственные средства
предприятия составляют 20 тыс, руб., сумма всех обязательств — 13 500 руб. Критерием
оценки выступает коэффициент риска. Сделать выводы.
\end{sloppypar}

Решение.

Коэф. риска K = Обязательства / собственные средства

K = 13.5 / 20 = 0.675

K < 1 - риска нет
K > 1 - высокий риск

Вывод: Риск банкротства низкий. Предприятие способно покрыть обязательства собственными средствами, финансовое положение устойчивое.


\section{Задание 10}
\begin{sloppypar}
Имеются два контракта на возведение одного и того же объекта в 300 млн руб. и
600 млн руб. Предполагается привлечение кредита на срок 4 года под 20\% годовых или
на срок 5 лет под 25\% годовых. Определить экономически выгодный вариант
строительства при ставке сравнения 18 %
\end{sloppypar}

\[
    a = \frac{i(1 + i)^n}{(1 + i)^n - 1}
\]

Вариант 1.

\[
    a_{4.20} = 0.385
\]

Готовые выплаты

\[
    A_1 = 300 * 0.385 = 115.5
\]

Приведенные затраты

\[
    Z_1 = 115.5 * \frac{ (1 - (1 + 0.18)^{-4} }{0.18} = 318.4 млн. руб.
\]


Вариант 2.


\[
    a_{5.25} = 0.358
\]

Готовые выплаты

\[
    A_1 = 300 * 0.385 = 115.5
\]

Приведенные затраты

\[
    Z_2 = 214.8 * \frac{ (1 - (1 + 0.18)^{-5} }{0.18} = 671.9 млн. руб.
\]

Вывод: Экономически выгоден первый вариант — строительство за 300 млн руб. с кредитом на 4 года под 20\%.

\newpage
Тестирование

\begin{itemize}
\item 1. б) Реорганизация юридического лица представляет собой прекращение юридического лица с переходом прав и обязанностей, то есть сохраняются права и обязанности при изменении организационно-правовой формы.
\item 2. а) Ликвидация юридического лица представляет собой прекращение деятельности без правопреемства, имущество и обязательства не переходят к другим лицам.
\item 3. в) Вступительный баланс составляется для оценки имущества на дату начала работы организации и отражает финансовое положение на стартовый момент.
\item 4. б) Ликвидационный баланс составляется для оценки имущества организации на момент ликвидации, чтобы правильно распределить активы между кредиторами и учредителями.
\item 5. а) При ликвидации задолженность по обязательным платежам в бюджет погашается вне установленной очереди, так как это первоочередные обязательства.
\item 6. а) Требования кредиторов по обязательствам, обеспеченным залогом имущества, удовлетворяются в первую очередь, так как залог гарантирует преимущественное право.
\item 7. а) Внутренние факторы, которые могут привести к банкротству, отслеживаются в ходе анализа отчета о финансовых результатах, так как он отражает прибыльность и финансовую устойчивость.
\item 8. б) Основными критериями неудовлетворительной структуры баланса неплатежеспособного предприятия являются коэффициент текущей ликвидности, коэффициент обеспеченности собственными средствами и коэффициент восстановления платежеспособности, так как они отражают способность предприятия покрывать свои обязательства.
\item 9. а) Внешнее управление назначается, когда платежеспособность предприятия может быть восстановлена, и целью является реабилитация предприятия.
\item 10. б) Конкурсный процесс — это специальная процедура ликвидации, обеспечивающая соразмерность требований кредиторов и охрану интересов всех сторон, контролируемая судом.
\item 11. б) Достижение договоренности между должником и кредитором относительно отсрочки или рассрочки платежей — это мировое соглашение, позволяющее частично урегулировать долговые обязательства без банкротства.
\item 12. б) Признание предприятия банкротом может быть осуществлено арбитражным судом, который официально устанавливает несостоятельность и назначает соответствующие процедуры.
\end{itemize}

\end{document}
